\chapter{优化}
\section{objreg索引优化}
//-----------------------------------------考虑优化--------------------
方案设计：

ObjReg是：
构建会维护一个模maxID 的递增计数器，为新对象分配唯一ID并且$$ID_i\in[0,maxID-1]$$以确保在最近$$2^k-1$$个区块内，该ID不重复。由于ObjReg是固定扇出为$$f$$的树，所以可将$$ID_i$$视为一个以$$f$$为基数的整数:
$$ID=\sum_{j=0}^{h-1}d_j\cdot F^j\quad\mathrm{where~}d_j\in[0,F-1]$$
每一位$$d_j$$对应ObjReg中第$$j$$层的子节点选择。由于Merkle树的子节点通常从1开始编号，因此树路径为$$tree\_index_{j} = d_j+1$$。如下图所示，假设ObjReg扇出$$f = 3$$，待插入对象$$o_6$$的ID为6，则$$6 = 0\times3^0+2\times3^1+0\times3^2$$，其三进制表示为020。因此其在ObjReg中的树路径为$$(1st,3rd,1st)$$,即第1，第2和第1个子节点。
在生成新区块时，我们除了插入新对象，还必须显示地将过期（超出滑动窗口）但尚未被覆盖的旧对象位置重置


ObjReg需要提取出数学原型，当前设计中所有对象都存储到叶子节点上了，大量的非叶节点仅仅用于哈希认证。这个索引设计效率太低。
并且，移除元素总是从最左侧元素移除，需要从底往上重新哈希。
原因在于假设位长为8，值为2，其二进制是00000010。我们依次找路径为1，1，1，1，1，1，2，1。如果我们认为前缀0无效，则其路径为2，1。倒过来是1，2.
现在对比：我们原来方案是认为二进制前导0都路径有效，所有元素存储在叶子节点。非叶节点仅仅计算其子节点哈希，然后层层向上计算出根哈希。
以扇出为3分析，依次插入第0-8个元素，位长为3. 二进制只能表示2位。所以我们才需要+1，因为根节点总是不存储任何信息。
ID/ 二进制/ 路径/   树中位置
0/  00/    11/    1层第1，2层第1，（这个第1第2仅仅针对该值在其父节点子孩子中的位置，1层也是我们将根节点视为第0层）
1/  01/    12/    1层第1，2层第2， 每个节点都会从根开始
2/  02/    13/    1层第1，2层第3.     
3/  10/    21/    1层第2，2层第1，
... 
8/  22/    33/    1层第3，2层第3. 
假设我们删除ID为0的元素，也就是往这个位置填入新的对象哈希，则依次会拿兄弟节点计算出根哈希。
举例：
树如下：h0
h1,h2,h3,
h4,h5,h6,h7,h8,h9,h10,h11,h12.
如果删除ID为0的对象，其所在节点哈希为h4，填充新对象，新对象会获得这个ID0，此时需要通过<h5,h6,h2,h3>向上重新计算根哈希。

如果扇出为偶，比如2，插入，只能插入0-7，完整树包括根所在的第0层，1-3层。所以扇出小，树就会深。但无论扇出奇偶我们计算路径总是需要在解析二进制后再+1
ID/ 二进制/ 路径/   树中位置
0/  000/    111/    1层第1，2层第1，3层1，（这个第1第2仅仅针对该值在其父节点子孩子中的位置，1层也是我们将根节点视为第0层）
1/  001/    12/    二层第1，三层第2， 每个节点都会从根开始
2/  010/    13/    根，二层第1，三层第3.     
3/  011/    21/    根，二层第2，三层第1，
... 
7/  111/    222/    1层第2，2层第2，3层第2。



现在方案认为前导0无效，也就是非叶节点会存储元素。
以扇出为3分析，依次插入第0-8个元素。
ID/ 二进制/ 路径/   树中位置
0/  00/    */    根。（我们认为*就是根，也就是前导全为0，移除之后仅剩*代表。）
1/  01/    1/    根，第1层第1
2/  02/    2/    根，第1层第2     
3/  10/    10/    
... 
我算不出这个树中位置了。我本意是，比如ID为1的元素，那我仅仅需要<h2,h3>做一层哈希。如果新方案能设计出来，那么我的非叶节点和叶节点都能存数据。会压缩树大小。
我不知道这新方案能不能实现。

我问了别人意见：原方案滑动窗口是模2^k-1个连续区间，ID按顺序递增，窗口整体向右移动。在新方案中，由于ID前缀不同，所以被移除的ID可能位于不同深度。若窗口大小固定为W，在任意时刻对象最多存储W个对象。
移除的ID与新增的ID可能享有共同前缀，因此树的而机构会动态变化？这个有点否认。
可能需要处理节点分裂与合并？我们好像只需要置空就可以？好像也不需要触发树结构调整。我有些想不明白。

其中有点不能否认,就是路径长度不一致，基于移动窗口方案，增加或者删除上可能不如原方案丝滑。目前我们有关所有的节点都用vec存储，实际上slots就存储底层那一排。这非常方便。先按照原定的来这个放到优化上去。


有关程序上不理解的：
1. 工程设计上暂时并没有用常见的LeafNode,Nonleafnode定义。
2. Result签名有问题。


代码设计：
1. 每次算根或proof，都需要把整棵树的各层哈希重新构建出来。性能很一般
2. 可以后期考虑改成增量上下文形式，升级：节点结构体，节点ID，写入上下文...


\section{关键字索引优化}
关键字索引构建

优化方式：
原有方式是按照ID决定插入位置，所以最终决定的树结构会和objreg一模一样，甚至可以合并为一个树。只不过我们为了设计更清晰，所以独立出一个树。
但是这样不能按照属性相似聚集在一起，以便给出更少的不相交证明。所以有什么更好的插入方式，更好的树组织方式能够减少给出不相交证明，聚合相似节点在一起。
底层叶子节点是一个对象及其属性。当插入一个新节点时，我们先将已经过期的元素删除，从上至下计算相似度决定往下走哪个分支，直至叶子节点。
父节点会存储子节点的所有属性，分为子节点相交属性，还有除相交之外属性（附加子节点映射）。举例来说，子节点属性是o1:a,b,c。o2:a,b,d。那么父节点就是:a,b,(c->o1),(d->o2).
查询时，先与所有节点验证相交，如果都没有，则折叠成哈希。如果都有，则将该非叶节点包括的对象ID全部纳入结果，如果有一部分，则往下继续深入。
但这有个问题，就是叠加了滑动窗口后，对旧有数据删除很复杂，以及哈希更新也很复杂。方案设计还要考虑acc值更新，验证数据生成。
这样组织有点像属性集合上的BTree树，但不知道理论和实际上是否可以实现。